\section{Verified accuracy of finite good aggregations}
\label{sec:formal-aggregation-accuracy}

For $u,v\in\R^r$, $r\geq2$, define
\[
 f(u,v)=(\|u\|^2-1,\ \|v\|^2-1,\ 1/2-u\transpose v),\qquad
 S=\{x:f_i(x)<0\text{ for all }i\},\qquad C=\cl(\conv S).
\]
The preceding verified classification identifies good multipliers exactly
as nonzero $\lambda\geq0$ with $\lambda_3^2\leq4\lambda_1\lambda_2$.
Write $f_\lambda=\sum_i\lambda_i f_i$. For a finite set $F$ of good
multipliers, put
\[
 P_F=\{x:f_\lambda(x)\leq0\text{ for every }\lambda\in F\},
 \qquad e_N=\inf_{|F|\leq N}d_H(C,P_F).
\]
Here $d_H$ is the extended Hausdorff distance in the Euclidean norm on all
$2r$ original coordinates. It is infinite for an unbounded relaxation.
The infimum ranges over arbitrary good multipliers, including interior
ones and arbitrary positive rescalings; no minimizing family is assumed.

The Lean package proves, for every $N\geq2$,
\[
 \frac{\sqrt2(\log2)^2}{1600N^2}
 \leq e_N\leq
 \frac{5\sqrt2\pi^2}{16(N-1)^2}.
\]
The constants are independent of $r$. A linear homeomorphism into an
explicit Euclidean space proves agreement with the original variables,
their quadratic residuals, and their actual closed convex hull. The
default maximum norm on a function space is not used in these error
bounds.

\paragraph{A uniform upper bound.}
The angle multipliers
\[
 \lambda(\theta)=(\cos^2\theta,\sin^2\theta,
                  2\sin\theta\cos\theta),\qquad 0\leq\theta\leq\pi/2,
\]
are good, and all their weak inequalities describe $C$. The formal proof
includes the coordinate endpoints and does not require the associated
two-by-two slack matrix to be PSD: only its nonnegative directions are
tested.

For a sample set with covering radius $\rho$, the angular quadratic
$h_x(\theta)=-f_{\lambda(\theta)}(x)$ is bounded below by $-5\rho^2$
throughout the interval whenever the sampled cuts and both coordinate
bounds hold. The proof minimizes this continuous function. At an interior
minimum, its derivative is zero; an exact rotation identity and
$|\sin t|\leq|t|$ give the bound. Endpoint minima are already nonnegative.
This is an elementary alternative to the paper's second-derivative
Taylor estimate.

With $a=5\rho^2$, contraction by $s=(1+2a)^{-1/2}$ repairs every point of
the relaxation into $C$. The Euclidean displacement is at most $\sqrt2a$.
The $N$ equally spaced angles have covering radius
$\pi/(4(N-1))$, giving the displayed upper constant. The estimate concerns
the entire relaxation, not only selected witnesses.

\paragraph{A finite-grid lower bound.}
The verified lower proof uses $N+1$ parameters
$\tau_j=1+j/N$, $j=0,\ldots,N$, and the rational perturbation
$\eta=1/(400N^2)$. It realizes the Gram matrices
\[
 \begin{pmatrix}
 1-1/(10\tau_j)&2/5-\eta\\
 2/5-\eta&1-\tau_j/10
 \end{pmatrix}
\]
by actual vectors in every required dimension. If one good multiplier
excludes witnesses at two parameters $\tau,\sigma\in[1,2]$, the cone
inequality gives
\[
 (\tau-\sigma)^2<16\bigl((1+10\eta)^2-1\bigr)<1/N^2.
\]
Thus a cut excludes at most one grid point, and at most $N$ cuts leave
one witness admitted. This finite pigeonhole argument needs no
interval-covering theorem; it includes zero-coordinate and interior
multipliers and imposes no bound on their ratios.

At an admitted witness, the good multiplier $(\tau,1/\tau,2)$ has value
$2\eta$. Its quadratic has Lipschitz constant at most $10$ on the product
of the two Euclidean unit balls, where both the witness and $C$ lie.
Every hull point therefore has distance at least $\eta/5$ from the
witness. This proves the stronger lower estimate
\[
 d_H(C,P_F)\geq\frac1{2000N^2}
 \geq\frac{\sqrt2(\log2)^2}{1600N^2}.
\]
The final comparison is also proved in Lean. The main paper improves
the first lower constant to $\sqrt2/2000$ using a sharper Lipschitz
estimate; that improvement is not part of this formal package.

\paragraph{Exact rational coefficients.}
For $m\geq1$, use
\[
 (m^2,j^2,2mj),\qquad(j^2,m^2,2mj),\qquad 0\leq j\leq m.
\]
The common middle ray is counted once, leaving exactly $2m+1$ distinct
cuts. They are good and include both coordinate rays. Their angles have
covering radius at most $1/(2m)$, proved using the Lipschitz bound for
arctangent. Positive rescaling connects these exact integer coefficients
to the angle-family theorem, yielding error at most
$5\sqrt2/(4m^2)$.

For $N\geq3$, choosing $m=\lfloor(N-1)/2\rfloor$ meets the cut budget.
Each coefficient is a nonnegative integer at most $2m^2$. With the binary
size convention assigning zero size to zero, the package bounds each size
by $2\operatorname{size}(N)+1$, and hence by
$2\lfloor\log_2 N\rfloor+3$. It also proves a uniform constant times
$N^{-2}$ error bound for these rational families. These statements concern
exact representation; they do not analyze numerical rounding or solver
conditioning.

\paragraph{Rates, tolerances, and scope.}
The extended-real family bounds imply finiteness of $e_N$ for $N\geq2$
before it is converted to a real number. The package then proves the
asymptotic $\Theta(N^{-2})$ statement. Writing
$c=\sqrt2(\log2)^2/1600$ and $B=5\sqrt2\pi^2/16$, any family of at most
$N\geq1$ cuts with error at most $\varepsilon>0$ must satisfy
$N\geq\sqrt{c/\varepsilon}$. Conversely the explicit angle family with at most
$\lceil\sqrt{B/\varepsilon}\rceil+1$ cuts suffices; its size is less than
$\sqrt{B/\varepsilon}+2$. Sufficiency constructs a family and does not
infer attainment from an infimum estimate.

The frozen claims and declaration coverage
are in \path{supplement/lean/audits/31-aggregation-accuracy/}. The support-function
interpretation and single-objective exact-aggregation proposition in the
main paper are outside this package. These uniform approximation
bounds give no solver runtime, iteration, or branch-and-bound lower bound,
and make no literature-priority claim.
